\chapter{Cross-Cloud Service Equivalence}
\index{cross-cloud equivalence}\index{service mapping}\index{multi-cloud}%
This appendix is a living reference that maps conceptually equivalent services and
ideas across the seven platforms studied in this book series: Databricks, Amazon Web
Services (AWS), Microsoft Azure, Microsoft Fabric, Google Cloud Platform (GCP),
MongoDB, and Snowflake. The names differ, but the underlying data-engineering
concepts---durable object storage, open table formats, elastic compute, declarative
pipelines, streaming, and disciplined governance---recur everywhere.

\begin{notebox}
This appendix consolidates the equivalence tables from every chapter's cross-cloud
callout---Chapters~1 through 24---into a single Rosetta Stone for moving fluently
between platforms. Where a clean one-to-one mapping does not exist, the prose notes the
closest analog and the caveat.
\end{notebox}

\section{Platform Foundations and Compute}
\index{compute!equivalence}\index{workspace!equivalence}%
Databricks' Chapter~1 topic---workspaces, clusters, and the control/data plane
split---maps to each cloud's way of bundling compute and administration.

\begin{center}
\footnotesize
\begin{tabular}{@{}p{3.0cm}p{3.4cm}p{3.2cm}p{3.2cm}@{}}
\toprule
\textbf{Concept} & \textbf{Databricks} & \textbf{AWS} & \textbf{Azure / GCP} \\
\midrule
Administrative unit & Workspace (+ account console) & Account / IAM Identity Center & Entra tenant / GCP project \\
Interactive compute & All-purpose cluster & EMR cluster / Glue interactive & Synapse Spark pool / Dataproc \\
Job compute & Job cluster (per run) & Glue job / EMR Serverless & Synapse job / Dataproc Serverless \\
Compute guardrails & Cluster policies + pools & IAM + service quotas & Policies / quotas \\
Billing unit & DBU + cloud VM cost & DPU / node-hours & vCore-hours / slot-hours \\
CLI / SDK & Databricks CLI, databricks-sdk-py & AWS CLI, Boto3 & az CLI / gcloud \\
\bottomrule
\end{tabular}
\end{center}

\noindent \textbf{Microsoft Fabric} replaces per-cluster decisions with a shared
capacity pool (F-SKU CUs); \textbf{Snowflake} replaces clusters with virtual
warehouses; \textbf{MongoDB Atlas} abstracts compute as cluster tiers.

\section{Table Storage and the Lakehouse}
\index{Delta Lake!equivalence}\index{lakehouse!equivalence}%
Databricks' Chapter~2 topic---Delta Lake---is the open table format at the center of
the lakehouse pattern on every cloud.

\begin{center}
\footnotesize
\begin{tabular}{@{}p{3.0cm}p{3.4cm}p{3.2cm}p{3.2cm}@{}}
\toprule
\textbf{Capability} & \textbf{Databricks} & \textbf{AWS / Azure / GCP} & \textbf{Snowflake / Fabric} \\
\midrule
Table format & Delta Lake & Delta / Iceberg / Hudi on object storage & Micro-partitions / Delta on OneLake \\
ACID on objects & Transaction log & Iceberg/Hudi catalogs; S3 Tables & Native / Delta transaction log \\
Time travel & \texttt{VERSION AS OF} & Iceberg snapshots & Time Travel / Delta versions \\
Compaction & \texttt{OPTIMIZE} & Iceberg \texttt{rewrite\_data\_files} & Automatic / \texttt{OPTIMIZE VORDER} \\
Clustering & Liquid clustering, Z-ORDER & Iceberg sort order, Athena partition & Clustering keys / V-Order + ZORDER \\
Row-level deletes & Deletion vectors & Iceberg position/m equality deletes & Native / deletion vectors \\
\bottomrule
\end{tabular}
\end{center}

\noindent \textbf{MongoDB} is the instructive contrast: a node-local document engine
(WiredTiger) rather than object-storage tables.

\section{The Medallion Pattern}
\index{medallion architecture!equivalence}%
Databricks' Chapter~3 topic---Bronze/Silver/Gold---is a portable design pattern
rather than a product.

\begin{center}
\footnotesize
\begin{tabular}{@{}p{2.6cm}p{3.4cm}p{3.4cm}p{3.0cm}@{}}
\toprule
\textbf{Layer} & \textbf{Databricks} & \textbf{Data lake (AWS/Azure/GCP)} & \textbf{Snowflake / Fabric} \\
\midrule
Bronze (raw) & \texttt{bronze} schema Delta tables & \texttt{raw/} prefix (S3/ADLS/GCS) & RAW schema / Bronze Lakehouse \\
Silver (validated) & \texttt{silver} schema Delta tables & \texttt{validated/} prefix & STAGING schema / Silver Lakehouse \\
Gold (curated) & \texttt{gold} schema star schema & \texttt{curated/} prefix & ANALYTICS schema / Gold Lakehouse \\
Governance & Unity Catalog grants per layer & Lake Formation / RBAC per zone & Role grants / workspace roles \\
\bottomrule
\end{tabular}
\end{center}

\section{Managed Spark and Execution Engines}
\index{Apache Spark!equivalence}\index{Photon!equivalence}%
Databricks' Chapter~4 topic---managed Apache Spark---is the reference implementation of
managed Spark; the engine is the same everywhere, but the operational wrapper differs.

\begin{center}
\footnotesize
\begin{tabular}{@{}p{3.0cm}p{3.4cm}p{3.2cm}p{3.2cm}@{}}
\toprule
\textbf{Capability} & \textbf{Databricks} & \textbf{AWS} & \textbf{Azure / GCP} \\
\midrule
Managed Spark & Databricks Runtime (Photon) & Glue / EMR / EMR Serverless & Synapse Spark / Fabric / Dataproc \\
Cluster manager & Built-in (no YARN to run) & YARN / EMR managed & Synapse / Dataproc managed \\
GPU/vectorized engine & Photon & EMR runtime (partial) & Fabric native engine \\
Autoscaling & Autoscaling clusters + serverless & EMR managed scaling & Autoscale pools / Dataproc autoscaling \\
UI and diagnostics & Spark UI + Ganglia + query profile & Spark UI on EMR & Spark UI / Fabric monitoring \\
Default table format & Delta & Iceberg/Hudi/Delta by config & Delta (Fabric) / BigQuery tables \\
\bottomrule
\end{tabular}
\end{center}

\noindent \textbf{Snowflake} plays this role with Snowpark (DataFrame API on its engine);
\textbf{MongoDB} integrates via the Spark Connector. Skills from Chapter~4 transfer to every
managed-Spark offering almost unchanged.

\section{SQL Warehousing and BI}
\index{SQL warehouse!equivalence}\index{BI!equivalence}%
Chapters~5--7 cover optimization, SQL warehouses, and AI/BI; every platform fields a
SQL serving layer and a BI surface.

\begin{center}
\footnotesize
\begin{tabular}{@{}p{3.0cm}p{3.4cm}p{3.2cm}p{3.2cm}@{}}
\toprule
\textbf{Capability} & \textbf{Databricks} & \textbf{AWS} & \textbf{Azure / GCP / Fabric} \\
\midrule
SQL compute & SQL warehouse (serverless/pro/classic) & Redshift Serverless / Athena & Synapse / BigQuery slots / Fabric Warehouse \\
Execution engine & Photon & Redshift engine / Athena (Trino) & Synapse / Dremel (BigQuery) / Fabric engine \\
Result reuse & Result cache + materialized views & Result cache / Athena reuse & Result-set cache / BI Engine / Direct Lake \\
BI surface & AI/BI dashboards + Genie & QuickSight (+ Q) & Power BI / Looker / Power BI in Fabric \\
NL analytics & Genie spaces & Amazon Q in QuickSight & Copilot / Gemini in Looker / Fabric Copilot \\
Cost introspection & \texttt{system.billing.usage} & CUR + CloudWatch & Cost Management / billing export / Capacity Metrics \\
\bottomrule
\end{tabular}
\end{center}

\section{Governance and Catalogs}
\index{governance!equivalence}\index{Unity Catalog!equivalence}%
Chapters~8 and 22 cover Unity Catalog; each hyperscaler answers governance
differently.

\begin{center}
\footnotesize
\begin{tabular}{@{}p{3.0cm}p{3.4cm}p{3.2cm}p{3.2cm}@{}}
\toprule
\textbf{Capability} & \textbf{Databricks} & \textbf{AWS} & \textbf{Azure / GCP} \\
\midrule
Catalog + governance & Unity Catalog & Glue Catalog + Lake Formation & Purview + OneLake / Dataplex + BigLake \\
Namespace & metastore $\rightarrow$ catalog $\rightarrow$ schema $\rightarrow$ table & account $\rightarrow$ database $\rightarrow$ table & tenant / project $\rightarrow$ dataset $\rightarrow$ table \\
Row/column security & Dynamic views, row filters, column masks & LF cell filters & RLS + OPA / policy tags \\
Lineage & UC lineage (table + column) & Glue lineage / SageMaker & Purview / Dataplex lineage \\
Audit & \texttt{system.access.audit} & CloudTrail + LF audit & Activity logs / audit logs \\
Secure sharing & Delta Sharing (open protocol) & S3 Access Grants / Redshift sharing & Azure Data Share / Analytics Hub \\
\bottomrule
\end{tabular}
\end{center}

\noindent \textbf{Snowflake} governs via account-level role hierarchies and shares
via Secure Data Sharing; \textbf{MongoDB} governs at the cluster/collection level
with Atlas roles and auditing.

\section{Declarative Pipelines, Ingestion, and Orchestration}
\index{pipelines!equivalence}\index{orchestration!equivalence}%
Chapters~9--12 cover PySpark tuning, DLT, Auto Loader, and Workflows.

\begin{center}
\footnotesize
\begin{tabular}{@{}p{3.0cm}p{3.4cm}p{3.2cm}p{3.2cm}@{}}
\toprule
\textbf{Capability} & \textbf{Databricks} & \textbf{AWS} & \textbf{Azure / GCP / Fabric} \\
\midrule
Declarative pipelines & Delta Live Tables (expectations) & Glue Data Quality + workflows & Fabric pipelines + DQ gates / Dataform \\
Incremental file ingest & Auto Loader (\texttt{cloudFiles}) & S3 events + Glue bookmarks & Event Grid + ADF / Storage transfer + Dataflow \\
CDC application & \texttt{MERGE} / \texttt{AUTO CDC} & DMS CDC to S3 + Glue & ADF CDC / Datastream + BigQuery MERGE \\
Orchestration & Workflows (jobs + tasks) & Step Functions / MWAA & ADF / Fabric pipelines / Cloud Composer \\
Pipeline as code & Asset Bundles + Terraform & CDK / Terraform & Bicep / Terraform / Fabric Terraform provider \\
\bottomrule
\end{tabular}
\end{center}

\section{Streaming}
\index{streaming!equivalence}%
Chapters~13--16 cover Structured Streaming and messaging integration.

\begin{center}
\footnotesize
\begin{tabular}{@{}p{3.0cm}p{3.4cm}p{3.2cm}p{3.2cm}@{}}
\toprule
\textbf{Capability} & \textbf{Databricks} & \textbf{AWS} & \textbf{Azure / GCP / Fabric} \\
\midrule
Stream processing & Structured Streaming & Flink (Managed Service) / Lambda & Stream Analytics / Dataflow / Fabric Eventstream \\
Message bus & Kafka (external), Event Hubs, Kinesis & Kinesis / MSK & Event Hubs / Pub/Sub / Eventstream \\
Exactly-once & Checkpoint + idempotent sink & Flink checkpoints + 2PC sinks & Dataflow exactly-once / at-least-once + dedup \\
State backend & RocksDB state store & RocksDB (Flink) & Managed state / in-memory \\
Real-time serving & Gold streaming tables + AI/BI & Kinesis Analytics $\rightarrow$ QuickSight & KQL DB + Real-Time Dashboards \\
\bottomrule
\end{tabular}
\end{center}

\section{Machine Learning and MLOps}
\index{MLOps!equivalence}\index{MLflow!equivalence}%
Chapters~17--20 cover MLflow, the Feature Store, Mosaic AI, and serving.

\begin{center}
\footnotesize
\begin{tabular}{@{}p{3.0cm}p{3.4cm}p{3.2cm}p{3.2cm}@{}}
\toprule
\textbf{Capability} & \textbf{Databricks} & \textbf{AWS} & \textbf{Azure / GCP / Fabric} \\
\midrule
Experiment tracking & MLflow (managed) & SageMaker Experiments & Azure ML / Vertex AI / Fabric MLflow \\
Model registry & Unity Catalog registry & SageMaker Model Registry & Azure ML registry / Vertex registry \\
Feature store & Feature Engineering in UC & SageMaker Feature Store & Azure ML FS / Vertex FS / Fabric feature sets \\
Vector search & Mosaic AI Vector Search & OpenSearch / Kendra vectors & AI Search / Vertex Vector Search / Fabric AI \\
GenAI platform & Mosaic AI (Foundation Models) & Bedrock & Azure OpenAI / Vertex AI / Fabric AI skills \\
Serving & Model Serving endpoints & SageMaker endpoints & Azure ML online / Vertex endpoints \\
Monitoring & Lakehouse Monitoring & SageMaker Model Monitor & Azure ML monitoring / Vertex monitoring \\
\bottomrule
\end{tabular}
\end{center}

\noindent \textbf{Snowflake} answers this layer with Snowpark ML, the Feature Store,
and Cortex; \textbf{MongoDB} with Atlas Vector Search and stream processing.

\section{Security, Identity, and Networking}
\index{security!equivalence}\index{identity!equivalence}\index{networking!equivalence}%
Databricks' Chapter~21 topic---identity, secrets, and network isolation---is a cloud-native
discipline with a Databricks control-plane twist.

\begin{center}
\footnotesize
\begin{tabular}{@{}p{3.0cm}p{3.4cm}p{3.2cm}p{3.2cm}@{}}
\toprule
\textbf{Capability} & \textbf{Databricks} & \textbf{AWS} & \textbf{Azure / GCP} \\
\midrule
Workforce identity & SSO + SCIM from IdP & IAM Identity Center & Entra ID / Cloud Identity \\
Machine identity & Service principal + OAuth M2M & IAM role & Managed identity / service account \\
Long-lived secret & PAT (policy-controlled) & Access keys & Client secrets / SA keys \\
Secrets store & Secret scopes (DB- or vault-backed) & Secrets Manager & Key Vault / Secret Manager \\
Private link & Private connectivity to control plane + storage & PrivateLink + VPC endpoints & Private Link / Private Service Connect \\
Perimeter control & IP access lists & Security groups, SCPs & NSGs, firewall rules / VPC SC \\
Audit & \texttt{system.access.audit} + auth logs & CloudTrail & Activity Log / Cloud Audit Logs \\
\bottomrule
\end{tabular}
\end{center}

\noindent \textbf{Snowflake} mirrors this with SSO/SCIM, network policies, and
Tri-Secret/CMK options; \textbf{MongoDB Atlas} with federated auth, private endpoints, and
IP access lists---the same checklist, different console.

\section{Infrastructure as Code and CI/CD}
\index{infrastructure as code!equivalence}\index{CI/CD!equivalence}\index{Asset Bundles!equivalence}%
Databricks' Chapter~23 topic---Asset Bundles, Terraform, and the delivery pipeline---has a
native answer per platform, while Terraform is the cross-cloud constant.

\begin{center}
\footnotesize
\begin{tabular}{@{}p{3.0cm}p{3.4cm}p{3.2cm}p{3.2cm}@{}}
\toprule
\textbf{Capability} & \textbf{Databricks} & \textbf{AWS} & \textbf{Azure / GCP / Fabric} \\
\midrule
Native IaC & Asset Bundles (\texttt{databricks.yml}) & CloudFormation / CDK & Bicep / Deployment Manager / Fabric Git + Terraform \\
Cross-platform IaC & Terraform provider & Terraform & Terraform \\
Validate-before-deploy & \texttt{bundle validate} & cfn-lint / CDK synth & \texttt{terraform plan} / what-if \\
CI auth & Service principal + OAuth M2M & OIDC $\rightarrow$ IAM role & OIDC $\rightarrow$ workload identity \\
State & Terraform remote state & S3 + DynamoDB lock & Storage account / GCS backend \\
Cost telemetry & \texttt{system.billing.usage} & CUR 2.0 & Cost Management / billing export / Capacity Metrics \\
\bottomrule
\end{tabular}
\end{center}

\noindent \textbf{Snowflake}'s story is Terraform-first plus database-change tools
(schemachange, dbt); the lesson that transfers everywhere: \emph{the merge is the deploy, and
the pipeline is the only writer to prod}.

\section{Certification and Readiness}
\index{certification!equivalence}%
Databricks' Chapter~24 topic---certification readiness---looks different per platform, but the
readiness method is identical: domains mapped to labs, mocks under time, remediation by
rebuilding.

\begin{center}
\footnotesize
\begin{tabular}{@{}p{3.0cm}p{3.4cm}p{3.2cm}p{3.2cm}@{}}
\toprule
\textbf{Element} & \textbf{Databricks} & \textbf{AWS} & \textbf{Azure / GCP / Fabric} \\
\midrule
Target exam & Data Engineer Associate / Professional & Data Engineer Associate & DP-700 / Professional Data Engineer / DP-600 \\
Heavy domains & Platform, Delta Lake, Spark ETL, pipelines, governance & Storage, ingestion, transform, ops & Fabric workloads / GCP services / Power BI \\
Lab weight & High (build fluency) & High & High \\
Mock standard & Two mocks, 80\%+, no domain $<$70\% & Same & Same \\
Capstone defense & Architecture review with cost numbers & Same & Same \\
\bottomrule
\end{tabular}
\end{center}

\noindent The portability argument of the whole series: if you can defend the Databricks
capstone \emph{and} sketch its AWS or Fabric equivalent from the cross-cloud boxes, you have
learned data engineering, not one console.

\begin{tipbox}
Use this appendix as a study aid for the book's lockstep, cross-cloud approach: when
you learn a concept on Databricks, immediately locate its equivalent here and
predict how the other platforms would express the same idea. That habit turns seven
separate vocabularies into one mental model.
\end{tipbox}
